\documentclass[11pt,a4paper]{article}
\usepackage[utf8]{inputenc}
\usepackage{amsmath,amssymb,amsthm}
\usepackage{listings}
\usepackage{xcolor}
\usepackage[margin=2.5cm]{geometry}
\usepackage{hyperref}

\lstset{
  language=C++,
  basicstyle=\ttfamily\small,
  keywordstyle=\color{blue}\bfseries,
  commentstyle=\color{green!60!black},
  stringstyle=\color{red!70!black},
  numbers=left,
  numberstyle=\tiny\color{gray},
  breaklines=true,
  frame=single,
  tabsize=4,
  showstringspaces=false
}

\title{IOI 2000 -- Post Office}
\author{Solution}
\date{}

\begin{document}
\maketitle

\section{Problem Statement}

There are $V$ villages on a line at sorted positions $x_1 < x_2 < \cdots < x_V$.
Place $P$ post offices at village locations to minimize the total distance from
each village to its nearest post office:
\[
  \min \sum_{i=1}^{V} \min_{j \in \text{post offices}} |x_i - x_j|.
\]

\noindent\textbf{Constraints:} $1 \le P \le V \le 300$.

\section{Solution Approach}

\subsection{Precomputation: Cost of Serving a Range}

For a contiguous group of villages $[i, j]$ served by a single optimally-placed post
office, the optimal location is the \textbf{median} village $m = \lfloor(i+j)/2\rfloor$.

Define:
\[
  w[i][j] = \sum_{k=i}^{j} |x_k - x_m|, \quad m = \lfloor(i+j)/2\rfloor.
\]

\begin{proof}[Proof of the recurrence $w[i][j] = w[i][j{-}1] + x_j - x_m$]
When extending from $[i, j{-}1]$ to $[i, j]$, the new median is
$m' = \lfloor(i+j)/2\rfloor$. If $j - i$ is even, then $m' = m + 1$ where
$m = \lfloor(i+j-1)/2\rfloor$ was the previous median. If $j - i$ is odd,
then $m' = m$. In either case, the incremental cost formula
$w[i][j] = w[i][j{-}1] + x_j - x_{m'}$ holds. This can be verified by
noting that moving the median from $m$ to $m'$ changes costs symmetrically for the
elements that switch sides, and $x_j - x_{m'}$ is the cost for the new village.
\end{proof}

We compute $w[i][j]$ incrementally in $O(V^2)$ total time.

\subsection{DP Formulation}

Let $\mathrm{dp}[i][j]$ = minimum total distance for villages $1, \ldots, i$ using
$j$ post offices.

\textbf{Base case:} $\mathrm{dp}[i][1] = w[1][i]$.

\textbf{Transition:}
\[
  \mathrm{dp}[i][j] = \min_{k = j-1}^{i-1} \bigl(\mathrm{dp}[k][j{-}1] + w[k{+}1][i]\bigr).
\]

\textbf{Answer:} $\mathrm{dp}[V][P]$.

\subsection{Reconstruction}

Store the optimal split point $k$ in an auxiliary array $\mathrm{opt}[i][j]$. Trace
backward to identify which group of villages each post office serves, placing each
post office at the median of its group.

\section{C++ Solution}

\begin{lstlisting}
#include <cstdio>
#include <algorithm>
#include <vector>
using namespace std;

int main() {
    int V, P;
    scanf("%d %d", &V, &P);

    vector<int> x(V + 1);
    for (int i = 1; i <= V; i++)
        scanf("%d", &x[i]);

    // Precompute w[i][j] = cost of one post office serving villages i..j
    vector<vector<int>> w(V + 1, vector<int>(V + 1, 0));
    for (int i = 1; i <= V; i++) {
        for (int j = i + 1; j <= V; j++) {
            int med = (i + j) / 2;
            w[i][j] = w[i][j - 1] + x[j] - x[med];
        }
    }

    // DP
    const int INF = 1000000000;
    vector<vector<int>> dp(V + 1, vector<int>(P + 1, INF));
    vector<vector<int>> opt(V + 1, vector<int>(P + 1, 0));

    for (int i = 1; i <= V; i++)
        dp[i][1] = w[1][i];

    for (int j = 2; j <= P; j++) {
        for (int i = j; i <= V; i++) {
            for (int k = j - 1; k <= i - 1; k++) {
                int cost = dp[k][j - 1] + w[k + 1][i];
                if (cost < dp[i][j]) {
                    dp[i][j] = cost;
                    opt[i][j] = k;
                }
            }
        }
    }

    printf("%d\n", dp[V][P]);

    // Reconstruct post office positions
    vector<int> postOffices;
    int curV = V, curP = P;
    while (curP > 1) {
        int k = opt[curV][curP];
        int med = (k + 1 + curV) / 2;
        postOffices.push_back(x[med]);
        curV = k;
        curP--;
    }
    int med = (1 + curV) / 2;
    postOffices.push_back(x[med]);

    reverse(postOffices.begin(), postOffices.end());
    for (int i = 0; i < (int)postOffices.size(); i++)
        printf("%d%c", postOffices[i], i + 1 < (int)postOffices.size() ? ' ' : '\n');

    return 0;
}
\end{lstlisting}

\section{Complexity Analysis}

\begin{itemize}
  \item \textbf{Time complexity:} $O(V^2 \cdot P)$ for the DP, plus $O(V^2)$ for
        precomputing $w$. With $V, P \le 300$: at most $300^2 \times 300 = 27{,}000{,}000$
        operations.
  \item \textbf{Space complexity:} $O(V^2 + V \cdot P)$.
\end{itemize}

\subsection{Possible Optimization}

The optimal split point $\mathrm{opt}[i][j]$ is non-decreasing in $i$ for fixed $j$
(the ``divide and conquer'' or ``Knuth'' optimization). This reduces the DP to
$O(V \cdot P)$ amortized. However, $O(V^2 P)$ is sufficient for $V \le 300$.

\end{document}
